\documentclass{article}
\usepackage{/home/user1/code/tex/sty}
\usepackage{amsmath}
\usepackage{amsfonts}
\usepackage{geometry}
\usepackage{hyperref}
\usepackage{floatrow}
\usepackage{array}
\newcolumntype{L}{>{$}l<{$}}
\newcommand{\tabb}[2]{
\begin{table}[H]
\begin{tabular}{#1}
#2
\end{tabular}
\end{table}}
\setlength{\parindent}{0pt}
\geometry{top=1in,bottom=1in,left=1in,right=1in}
\begin{document}
\begin{center}
\textbf{\Large{Problem Set 6 Corrections}}\\~\\
\begin{tabular}{l l}
Student&Agustin Romero\\
Email&ar1@stanford.edu\\
Instructor&Caterina Vernieri\\
Course&Physics 252\\
Institution&Stanford University\\
Date&\today\\
\end{tabular}
\end{center}
\section*{Problem 1}
The student answer was not correct, because it did not include the list of probabilities. The corrected solution is below.\\
\hrule
~\\~\\
From the Clebsch-Gordon table,
\e{\ket{2,0}\times \ket{2,0}=(\fr{18}{35})^{1/2}\ket{4,0}+0\ket{3,0}-(\fr{2}{7})^{1/2}\ket{2,0}+0\ket{1,0}+(\fr{1}{5})^{1/2}\ket{0,0}}
The probability of each total angular momentum state $j$ is
\tabb{L L}{
j & \textrm{Probability}\\
0&\fr{1}{5}\\
1&0\\
2&\fr{2}{7}\\
3&0\\
4&\fr{18}{35}}
The sum of the probabilities is
\e{\sum P = \fr{1}{5}+\fr{2}{7}+\fr{18}{35}=\boxed{1}}
\section*{Problem 2}
The student solution was correct.
\section*{Problem 3.a.}
The student solution was almost correct, except it should not have the additional permutation of 2 S states, and the state should be multiplied out.\\
\hrule
~\\~\\
The antisymmetric baryonic wavefunction for $\Omega^{-}$ with all spins aligned is
\e{\psi=\psi_{Ss}\psi_{Sf}\psi_{Ac}}
The antisymmetric baryonic color wavefunction is
\e{\psi_{Ac}=\fr{1}{\sqrt{6}}(rgb-rbg+gbr-grb+brg-bgr)}
The symmetric spin wavefunction $\ket{\fr{3}{2},-\fr{3}{2}}$ is
\e{\psi_{Ss}=(\downarrow_1 \downarrow_2 \downarrow_3)}
The symmetric flavor wavefunction for $\Omega^{-}$ is
\e{\psi_{Sf}=sss}
Finally the wavefunction for $\Omega^{-}$ with $m_s=-\fr{3}{2}$ is
\[ \! \boxed{\begin{aligned}
\psi&=\fr{1}{\sqrt{6}}(\downarrow_1 \downarrow_2 \downarrow_3)(sss)(rgb-rbg+gbr-grb+brg-bgr)\\
&= \fr{1}{\sqrt{6}}((s\downarrow r)(s\downarrow g)(s\downarrow b)-(s\downarrow g)(s\downarrow r)(s\downarrow b)+(s\downarrow g)(s\downarrow b)(s\downarrow r)\\
&-(s\downarrow b)(s\downarrow g)(s\downarrow r)+(s\downarrow b)(s\downarrow r)(s\downarrow g)-(s\downarrow r)(s\downarrow b)(s\downarrow g))
\end{aligned}} \]
\section*{Problem 3.b.}
The student solution was correct.
\section*{Problem 4.}
The student solution was not correct, because the derived multiplicities were not correct.\\
\hrule
~\\~\\
For mesons, there are no spin degrees of freedom, and there are 3 available color states, and 9 available flavor states. The parity of the meson is
\e{P=P_1 P_2 (-1)^L=(-1)^L}
and total angular momentum $L=0,1,...$. So the multiplicity of the $L=0$ state is $\boxed{9}$. There is one nonet associated with
\e{\boxed{J^{P}=0^+, 1^{-}, 2^{+}, ...}}
For baryons, the wavefunction is
\e{\psi=\psi_{\rom{flavor}}\psi_{\rom{color}}\psi_{\rom{space}}}
$\psi_{\rom{color}}$ is antisymmetric, then $\psi_{\rom{flavor}}\psi_{\rom{space}}$ is antisymmetric. When $L=0$, $\psi_{\rom{space}}$ is symmetric, therefore $\psi_{\rom{flavor}}$ is antisymmetric. Using the lightest quarks $u$, $d$, and $s$, the flavor wavefunction is
\e{\psi_{\rom{flavor}}=\fr{1}{\sqrt{6}}(uds-dus+dsu-sdu+sud-usd)}
The multiplicity of the baryonic state is $\boxed{1}$ with $J=0$, $L=0$, $P=(+1)^3(-1)^0=+1$.
\section*{Problem 5.a.}
The student solution was correct.
\section*{Problem 5.b.}
The student solution was not correct because $f$ was incorrectly calcualted.\\
\hrule
~\\~\\
The r, g, b states are
\e{r=\pma{1\\ 0 \\ 0} \quad b=\pma{0 \\ 1 \\ 0} \quad g=\pma{0 \\ 0 \\ 1}}
The color singlet state is
\e{\ket{8}=\fr{1}{\sqrt{3}}(r\bar{r}+b\bar{b}+g\bar{g})}
Then
\e{c_2^\dagger \lambda^\alpha c_1 &= \fr{1}{\sqrt{3}}(\bar{r}\lambda^\alpha r + \bar{b}\lambda^\alpha b + \bar{g}\lambda^\alpha g)\\
&= \fr{1}{\sqrt{3}}(\pma{1 & 0 & 0} \lambda^\alpha \pma{1 \\ 0 \\ 0} + \pma{0 & 1 & 0}\lambda^\alpha \pma{0 \\ 1 \\ 0} + \pma{0 & 0 & 1}\lambda^\alpha \pma{0 \\ 0 \\ 1})\\
&= \fr{1}{\sqrt{3}}(\lambda_{11}^\alpha+\lambda_{22}^\alpha+\lambda_{33}^\alpha)}
Then calculate $f$
\e{f=\fr{1}{4}\sum_{\alpha}(c_2^\dagger \lambda^\alpha c_1)^2 = \fr{1}{12}(1+\fr{1}{\sqrt{3}}-1+\fr{1}{\sqrt{3}}-\fr{2}{\sqrt{3}})^2=\boxed{0}}
The color factor $f$ is zero for this process assuming the quarks are in a color singlet state. Because the quark and antiquark are initially a color singlet state, the gluon \fbox{propagator would be colorless and noninteracting}.
\section*{Problem 6.}
The student solution was correct.
\end{document}
